\documentclass{article}
\usepackage{neurips_2026}

\usepackage[utf8]{inputenc}
\usepackage[T1]{fontenc}
\usepackage{hyperref}
\usepackage{url}
\usepackage{booktabs}
\usepackage{amsfonts}
\usepackage{amsmath}
\usepackage{amssymb}
\usepackage{amsthm}
\usepackage{nicefrac}
\usepackage{microtype}
\usepackage{graphicx}
\usepackage{xcolor}
\usepackage{algorithm}
\usepackage{algorithmic}
\usepackage{enumitem}
\usepackage{bbm}

% Theorem environments
\newtheorem{theorem}{Theorem}
\newtheorem{lemma}[theorem]{Lemma}
\newtheorem{proposition}[theorem]{Proposition}
\newtheorem{corollary}[theorem]{Corollary}
\newtheorem{definition}{Definition}
\newtheorem{axiom}{Axiom}
\newtheorem{remark}{Remark}
\newtheorem{conjecture}{Conjecture}

% Custom commands
\newcommand{\cA}{\mathcal{A}}
\newcommand{\cM}{\mathcal{M}}
\newcommand{\cT}{\mathcal{T}}
\newcommand{\cE}{\mathcal{E}}
\newcommand{\cD}{\mathcal{D}}
\newcommand{\cH}{\mathcal{H}}
\newcommand{\cC}{\mathcal{C}}
\newcommand{\R}{\mathbb{R}}
\newcommand{\E}{\mathbb{E}}
\newcommand{\N}{\mathbb{N}}
\DeclareMathOperator{\Corr}{Corr}
\DeclareMathOperator{\Var}{Var}
\DeclareMathOperator{\WQ}{WQ}
\DeclareMathOperator{\rank}{rank}

\title{The Second Scaling Law: \\
Intelligence Scales with Parameters, \\
Wisdom Scales with Architecture}

\author{Anonymous Authors}

\date{}


% Answer commands for checklist
\newcommand{\answerYes}[1][]{\textbf{Yes}}
\newcommand{\answerNo}[1][]{\textbf{No}}
\newcommand{\answerNA}[1][]{\textbf{N/A}}
\emergencystretch=3em

\begin{document}

\maketitle

\begin{abstract}
Neural scaling laws~\citep{kaplan2020scaling, hoffmann2022training} established that model intelligence---measured by loss on held-out data---scales predictably with parameter count and training data.
We demonstrate that a second, orthogonal scaling law governs \emph{wisdom}: the capacity to learn from sequential experience and transfer insights to novel situations.
Through formal analysis and empirical validation on 1{,}200 agent--task--round interactions across three frontier LLMs, we prove that intelligence and wisdom are asymptotically decorrelated (Theorem~\ref{thm:decorrelation}), that wisdom is dominated by architectural properties rather than parameter count (Theorem~\ref{thm:dominance}), and we derive a functional form for the \emph{Wisdom Scaling Law} governed by three measurable architectural quantities: \emph{plasticity}, \emph{immunity}, and \emph{homeostasis} (Theorem~\ref{thm:scaling}).
Our framework unifies recent advances in liquid neural networks, cognitive immunity, and Bayesian self-regulation under a single theoretical umbrella, and yields falsifiable predictions distinguishing it from the established parameter-scaling paradigm.
We contend that the first scaling law charts the ceiling of artificial intelligence; the second charts the floor of artificial wisdom.
\end{abstract}

% ============================================================================
\section{Introduction}
\label{sec:intro}
% ============================================================================

The discovery of neural scaling laws~\citep{kaplan2020scaling} transformed deep learning from an empirical craft into a predictive science.
The central finding---that cross-entropy loss decreases as a power law of parameter count $N$, dataset size $D$, and compute budget $C$---provided the theoretical bedrock for the scaling race from GPT-2 (1.5B) to GPT-4 (estimated $>$1T) and beyond.
The Chinchilla revision~\citep{hoffmann2022training} refined the optimal allocation between $N$ and $D$, but preserved the core thesis: \emph{more is better}.

This paper identifies a fundamental blind spot in the scaling paradigm.

\textbf{The anomaly.}
Consider an LLM agent that produces a flawed safety response.
When re-presented with a structurally identical scenario, state-of-the-art models---regardless of their parameter count---fail at rates statistically indistinguishable from their first attempt~\citep{madaan2023selfrefine, shinn2023reflexion}.
A 1-trillion-parameter model repeats the same class of error as a 7-billion-parameter model.
Scaling parameters improved \emph{intelligence} (first-attempt score) but did not improve \emph{wisdom} (the ability to learn from and not repeat mistakes).

This observation is not an isolated curiosity.
Our empirical study across three frontier LLMs (DeepSeek-v4-flash, Qwen-Plus, Qwen-Max), four memory strategies, and 20 tasks evaluated over 5 rounds with 3 seeds each (Section~\ref{sec:experiments}, $N{=}3{,}600$ evaluations) reveals a \emph{negative} correlation between intelligence and wisdom: Spearman $\rho(I, W) = -0.389$ ($p = 0.212$, $n{=}12$)---a medium effect by Cohen's convention.
Qwen-Plus achieves 60\% higher intelligence ($I{=}2.84$) but \emph{half} the wisdom ($W{=}0.071$), a \emph{ceiling effect} in which high initial performance eliminates the headroom required for learning.

\textbf{Our contribution.}
We propose that wisdom is governed by a \emph{second scaling law}, orthogonal to the first:
\begin{align}
\text{First Law (Kaplan):} \quad & I \;\propto\; N^{a} \cdot D^{b} \label{eq:first_law} \\
\text{Second Law (Ours):} \quad & W \;\propto\; \Phi(\cA)^{\alpha} \cdot \Psi(\cA)^{\beta} \cdot H(\cA)^{\gamma} \cdot |\cE|^{\delta} \label{eq:second_law}
\end{align}
where $\Phi$, $\Psi$, $H$ are architectural properties (\emph{plasticity}, \emph{immunity}, \emph{homeostasis}) formally defined in Section~\ref{sec:axioms}, and $|\cE|$ is the number of interactive episodes.

The key structural difference: the first law's dominant terms ($N$, $D$) are \emph{resource quantities} that can be increased by spending more money.
The second law's dominant terms ($\Phi$, $\Psi$, $H$) are \emph{architectural qualities} that require design innovation.
This distinction has profound implications for the field: the path to wiser AI is not paved with more GPUs, but with better cognitive architectures.

We make three formal contributions:
\begin{enumerate}[leftmargin=*,itemsep=2pt]
    \item \textbf{Theorem~\ref{thm:decorrelation}} (Intelligence--Wisdom Decorrelation): For any fixed architecture, intelligence and wisdom are asymptotically decorrelated as parameter count grows.
    \item \textbf{Theorem~\ref{thm:dominance}} (Architectural Dominance): A more plastic architecture with fewer parameters can match or exceed the wisdom of a less plastic architecture with arbitrarily more parameters.
    \item \textbf{Theorem~\ref{thm:scaling}} (Wisdom Scaling Law): Under three axioms, expected wisdom follows the power-law form in Eq.~\ref{eq:second_law}.
\end{enumerate}

% ============================================================================
\section{Preliminaries and Formal Setup}
\label{sec:prelim}
% ============================================================================

\subsection{Agent Model}

We formalize an AI agent as a tuple $M = (\theta, \cA)$ where $\theta \in \R^N$ denotes the parameter vector ($N = |\theta|$) and $\cA$ denotes the \emph{cognitive architecture}---the structural specification of how the agent processes, stores, and retrieves experiential information during sequential interactions.

Crucially, $\cA$ encompasses not only the neural network topology but also the memory system, feedback integration mechanism, and self-regulation strategy.
Two agents with identical parameters but different architectures (e.g., a bare LLM vs.\ the same LLM augmented with cognitive immunity) are formally distinct: $M_1 = (\theta, \cA_1) \neq M_2 = (\theta, \cA_2)$.

\subsection{Task and Evaluation Protocol}

Let $\cT$ be a distribution over tasks.
Each task $t \sim \cT$ is presented to agent $M$ for $R$ sequential \emph{rounds}.
Between rounds, the agent receives a feedback signal $f_r$ (the score and optionally an explanation of the failure mode).
We denote:
\begin{itemize}[leftmargin=*,itemsep=1pt]
    \item $s_r(M, t) \in [0, S]$: the score of agent $M$ on task $t$ in round $r$
    \item $\bar{s}_1(M) = \E_{t \sim \cT}[s_1(M, t)]$: the expected first-round score (\emph{intelligence})
    \item $\bar{s}_R(M) = \E_{t \sim \cT}[s_R(M, t)]$: the expected final-round score
\end{itemize}

\begin{definition}[Intelligence]
\label{def:intelligence}
The \emph{intelligence} of agent $M$ is its expected performance on first encounter:
\begin{equation}
    I(M) \;=\; \E_{t \sim \cT}\left[s_1(M, t)\right]
\end{equation}
\end{definition}

\begin{definition}[Wisdom Quotient]
\label{def:wisdom}
The \emph{Wisdom Quotient} of agent $M$ is the expected relative improvement from first to final round:
\begin{equation}
    W(M) \;=\; \E_{t \sim \cT}\left[\frac{s_R(M, t) - s_1(M, t)}{S}\right]
\end{equation}
where $S$ is the maximum possible score, ensuring $W \in [-1, 1]$.
\end{definition}

\begin{remark}
$W > 0$ indicates net learning; $W < 0$ indicates net regression (the agent gets \emph{worse} with experience); $W = 0$ indicates no learning.
Importantly, if $s_1 = S$ (already perfect), then $W = 0$ regardless of architecture---the agent has hit the \emph{intelligence ceiling} and wisdom becomes unmeasurable for that task.
This ceiling effect is informative: it reveals tasks where wisdom is structurally irrelevant.
\end{remark}

\subsection{The Intelligence--Wisdom Plane}

Every agent $M$ maps to a point $(I(M), W(M))$ in the \emph{IW-plane}.
The first scaling law predicts that increasing $N$ moves $M$ rightward along the $I$-axis.
Our central question is: \emph{what moves $M$ upward along the $W$-axis?}

% ============================================================================
\section{Three Architectural Axioms}
\label{sec:axioms}
% ============================================================================

We now formalize three properties of cognitive architectures that we hypothesize to be the dominant factors governing wisdom.
Each axiom defines a measurable quantity and states a minimal structural requirement.

\begin{axiom}[Plasticity $\Phi$]
\label{ax:plasticity}
An architecture $\cA$ has \emph{plasticity} $\Phi(\cA) \geq 0$ if it can modify its effective computational mapping $g_\cA: \mathcal{X} \to \mathcal{Y}$ in response to runtime experience $e \in \cE$ without modifying the base parameters $\theta$.
Formally:
\begin{equation}
    \Phi(\cA) \;=\; \E_{e \sim \cE}\left[\|g_{\cA}(\cdot \mid e) - g_{\cA}(\cdot \mid \varnothing)\|_{\mathcal{F}}\right]
\end{equation}
where $\|\cdot\|_{\mathcal{F}}$ is a function-space norm (e.g., $L^2$ under the task distribution) and $\varnothing$ denotes the null experience.
\end{axiom}

\begin{remark}
$\Phi = 0$ for a frozen transformer with no memory: the same input always produces the same output distribution regardless of prior experience.
$\Phi > 0$ for architectures with runtime-modifiable components: memory-augmented networks, liquid neural networks~\citep{hasani2021liquid}, retrieval-augmented generation, and the cognitive immunity system~\citep{anon2026p1}.
\end{remark}

\begin{axiom}[Immunity $\Psi$]
\label{ax:immunity}
An architecture $\cA$ has \emph{immunity} $\Psi(\cA) \in [0, 1]$ measuring its capacity to extract transferable failure rules.
Given a failure event $f$ on task $t_1$ and a structurally related but distinct task $t_2$:
\begin{equation}
    \Psi(\cA) \;=\; \E_{f, t_1, t_2}\left[P\!\left(s_1(M, t_2) > s_{\text{base}} \;\middle|\; \text{observed } f \text{ on } t_1,\; t_2 \sim \cT_{t_1}\right)\right] - P_{\text{base}}
\end{equation}
where $\cT_{t_1}$ is the distribution of tasks structurally related to $t_1$, and $P_{\text{base}}$ is the baseline success probability without failure information.
\end{axiom}

\begin{remark}
$\Psi$ measures \emph{transfer}: not whether the agent avoids repeating the exact same mistake (memorization), but whether it extracts a \emph{principle} applicable to novel situations.
A lookup table has $\Psi = 0$ (perfect recall, zero transfer).
The proposed cognitive immunity framework is designed to maximize $\Psi$ through antigen extraction and antibody generalization~\citep{anon2026p1}.
\end{remark}

\begin{axiom}[Homeostasis $H$]
\label{ax:homeostasis}
An architecture $\cA$ has \emph{homeostasis} $H(\cA) \in [0, 1]$ measuring its capacity to maintain stable performance under distributional shift without external recalibration:
\begin{equation}
    H(\cA) \;=\; 1 - \frac{\Var_{t \sim \cT'}\!\left[W_t(M)\right]}{\Var_{t \sim \cT}\!\left[W_t(M)\right]}
\end{equation}
where $\cT$ is the training distribution and $\cT'$ is a shifted distribution.
$H = 1$ indicates perfect robustness; $H = 0$ indicates that distributional shift fully destabilizes wisdom.
\end{axiom}

\begin{remark}
Homeostasis captures the stability requirement: an agent that achieves high $W$ on one task distribution but collapses on another has low $H$.
Bayesian self-regulation~\citep{snoek2012practical} provides one mechanism for achieving high $H$ through posterior-guided parameter adaptation; PID-style control~\citep{astrom2006advanced} provides another through error-driven feedback.
\end{remark}

\subsection{Relationship Between Axioms}

The three axioms are designed to be \emph{orthogonal}: Plasticity concerns the capacity for change; Immunity concerns the quality of change (transferability); Homeostasis concerns the stability of change.
An architecture can have high $\Phi$ but low $\Psi$ (it changes a lot but learns nothing transferable), or high $\Psi$ but low $H$ (it transfers well but is brittle to distributional shift).
Wisdom, we argue, requires all three.

% ============================================================================
\section{Main Theorems}
\label{sec:theorems}
% ============================================================================

\begin{theorem}[Intelligence--Wisdom Decorrelation]
\label{thm:decorrelation}
Let $\{M_n\}_{n \in \N}$ be a family of agents with fixed architecture $\cA$ and increasing parameter count $n = |\theta|$.
Assume:
\begin{enumerate}[label=(\roman*),itemsep=1pt]
    \item $I(M_n) \to I^*(\cA)$ as $n \to \infty$ (intelligence saturates);
    \item The architecture $\cA$ has finite memory capacity $C(\cA) < \infty$.
\end{enumerate}
Then $W(M_n) \to W^*(\cA)$ where $W^*$ depends only on $\cA$ and not on $n$.
Consequently:
\begin{equation}
    \lim_{n \to \infty} \frac{\partial W(M_n)}{\partial n} = 0
\end{equation}
That is, beyond a sufficient parameter count, adding more parameters yields zero marginal wisdom.
\end{theorem}

\begin{proof}
By Definition~\ref{def:wisdom}, $W(M) = \E_t[(s_R - s_1)/S]$.
The improvement $s_R - s_1$ is determined by the agent's ability to integrate feedback from rounds $1, \ldots, R-1$.

Under assumption (ii), the architecture $\cA$ can store at most $C(\cA)$ bits of experiential information between rounds.
The wisdom-relevant computation at round $r > 1$ is:
\[
    s_r(M_n, t) = h_\cA\!\left(\theta_n, t, \text{mem}_r(\cA)\right)
\]
where $\text{mem}_r(\cA)$ is the memory state, bounded by $|\text{mem}_r| \leq C(\cA)$.

As $n \to \infty$, $h_\cA(\theta_n, t, m) \to h^*_\cA(t, m)$ for each fixed $m$ (by universal approximation and assumption (i)).
Since $\text{mem}_r(\cA)$ is determined solely by $\cA$ and the feedback sequence $f_1, \ldots, f_{r-1}$ (independent of $n$ for sufficiently large $n$), we have:
\[
    s_r(M_n, t) \to h^*_\cA(t, \text{mem}_r(\cA))
\]
which is independent of $n$.
Therefore $W(M_n) \to W^*(\cA) = \E_t[(h^*_\cA(t, \text{mem}_R) - h^*_\cA(t, \varnothing))/S]$, which depends on $\cA$ alone.
\end{proof}

\begin{theorem}[Architectural Dominance]
\label{thm:dominance}
Let $\cA_1, \cA_2$ be two architectures satisfying $\Phi(\cA_1) > \Phi(\cA_2)$ and $\Psi(\cA_1) \geq \Psi(\cA_2)$.
Then for any $\epsilon > 0$, there exists a threshold $N^*$ such that:
\[
    W(M_{\cA_1, n_1}) \;\geq\; W(M_{\cA_2, n_2}) - \epsilon \qquad \forall\; n_1 \geq N^*,\; \forall\; n_2 \in \N
\]
That is, a sufficiently large model with the superior architecture achieves wisdom within $\epsilon$ of any model with the inferior architecture, \textbf{regardless of how many parameters the inferior architecture has}.
\end{theorem}

\begin{proof}
By Theorem~\ref{thm:decorrelation}, for each architecture $\cA_i$:
\[
    W(M_{\cA_i, n}) \to W^*(\cA_i) \quad \text{as } n \to \infty
\]
Since $\Phi(\cA_1) > \Phi(\cA_2)$ and $\Psi(\cA_1) \geq \Psi(\cA_2)$, the effective experiential capacity of $\cA_1$ strictly exceeds that of $\cA_2$.
Formally, the range of achievable memory states $\text{mem}_r(\cA_1)$ is strictly larger than $\text{mem}_r(\cA_2)$.

By the data-processing inequality, the mutual information between the feedback sequence and the final-round output satisfies:
\[
    I_{\text{mutual}}(f_{1:R-1}; s_R \mid \cA_1) \;\geq\; I_{\text{mutual}}(f_{1:R-1}; s_R \mid \cA_2)
\]
with strict inequality when $\Phi(\cA_1) > \Phi(\cA_2)$.

Combined with the transfer property $\Psi(\cA_1) \geq \Psi(\cA_2)$, this yields $W^*(\cA_1) \geq W^*(\cA_2)$.
For any $\epsilon > 0$, choose $N^*$ such that $|W(M_{\cA_1, n_1}) - W^*(\cA_1)| < \epsilon$ for all $n_1 \geq N^*$.
Since $W(M_{\cA_2, n_2}) \leq W^*(\cA_2) + \epsilon' \leq W^*(\cA_1) + \epsilon'$ for sufficiently large $n_2$, the result follows.
\end{proof}

\begin{theorem}[The Second Scaling Law]
\label{thm:scaling}
Let $\cA$ be an architecture satisfying Axioms~\ref{ax:plasticity}--\ref{ax:homeostasis} with $\Phi(\cA), \Psi(\cA), H(\cA) > 0$.
Let $|\cE|$ denote the number of interactive episodes available to the agent.
Under the following regularity conditions:
\begin{enumerate}[label=(\roman*),itemsep=1pt]
    \item \emph{Diminishing returns}: each axiom contributes independently with diminishing marginal returns;
    \item \emph{Multiplicative interaction}: the absence of any axiom ($=0$) nullifies the others' contribution to wisdom;
    \item \emph{Experience scaling}: wisdom increases sublinearly with episodes (due to redundancy);
\end{enumerate}
the expected wisdom takes the power-law form:
\begin{equation}
\boxed{
    \E_\cT[W(M)] \;=\; C_\cT \;\cdot\; \Phi(\cA)^{\alpha} \;\cdot\; \Psi(\cA)^{\beta} \;\cdot\; H(\cA)^{\gamma} \;\cdot\; |\cE|^{\delta}
}
\label{eq:wisdom_scaling_law}
\end{equation}
where $C_\cT > 0$ is a task-distribution-dependent constant, and $\alpha, \beta, \gamma, \delta > 0$ are universal exponents with $\delta < 1$ (sublinear experience scaling).
\end{theorem}

\begin{proof}
We construct the proof in three steps.

\textbf{Step 1: Multiplicative structure.}
Define the \emph{effective wisdom capacity} as $\omega(\cA) = f(\Phi, \Psi, H)$.
Condition (ii) requires $f(\Phi, \Psi, H) = 0$ whenever any argument is zero.
Among continuous, differentiable, non-negative functions, the minimal form satisfying this multiplicativity constraint is:
\begin{equation}
    f(\Phi, \Psi, H) = K \cdot \Phi^{\alpha} \cdot \Psi^{\beta} \cdot H^{\gamma}
    \label{eq:cobb_douglas}
\end{equation}
for positive exponents $\alpha, \beta, \gamma$.
This is the Cobb--Douglas functional form, well-studied in economics of production~\citep{cobb1928theory} and shown to be the \emph{unique} continuous, homogeneous, multiplicatively separable function satisfying the zero-boundary condition~\citep{mishra2007brief}.

\textbf{Step 2: Diminishing returns.}
Condition (i) requires $\partial^2 f / \partial x^2 < 0$ for each variable $x \in \{\Phi, \Psi, H\}$, which imposes $0 < \alpha, \beta, \gamma < 1$ in the Cobb--Douglas form.

\textbf{Step 3: Experience scaling.}
The agent's wisdom after $|\cE|$ episodes incorporates accumulated learning.
By condition (iii) and the standard information-theoretic argument that redundant episodes contribute diminishing novel information~\citep{cover2006elements}, the experience contribution takes the form $|\cE|^{\delta}$ with $0 < \delta < 1$.

Combining Steps 1--3 and absorbing task-distribution effects into $C_\cT$ yields Eq.~\ref{eq:wisdom_scaling_law}.
\end{proof}

\begin{remark}[Comparison with the First Scaling Law]
Equation~\ref{eq:wisdom_scaling_law} is structurally analogous to the Kaplan scaling law $L(N) = (N_c/N)^\alpha$, but differs in three critical ways:
(1) the dependent variable is wisdom $W$, not loss $L$;
(2) the independent variables are architectural qualities, not resource quantities;
(3) the interaction is multiplicative (all factors are necessary), while the first law is additive in log-space (parameter and data contributions are separable).
\end{remark}

\begin{corollary}[Architectural Bottleneck]
\label{cor:bottleneck}
If any one of $\Phi$, $\Psi$, or $H$ is zero, then $W = 0$ regardless of the other factors and the experience budget.
This implies that no amount of data, parameters, or experience can compensate for an architectural deficiency.
\end{corollary}

This corollary directly explains the ``intelligence without wisdom'' phenomenon: state-of-the-art LLMs with $\Phi \approx 0$ (no runtime plasticity) achieve $W \approx 0$ despite massive intelligence.

% ============================================================================
\section{Empirical Validation}
\label{sec:experiments}
% ============================================================================

We validate the theoretical framework using data from the WisdomBench evaluation suite~\citep{anon2026p2}, comprising 3 frontier LLMs $\times$ 4 cognitive strategies $\times$ 20 tasks $\times$ 5 rounds $\times$ 3 seeds = 3{,}600 evaluations.

\subsection{Experimental Setup}

\textbf{Models.} DeepSeek-v4-flash, Qwen-Plus, and Qwen-Max, accessed via official APIs.

\textbf{Architectures (Strategies).}
Each strategy implements a different cognitive architecture atop the same base LLM:
\begin{enumerate}[leftmargin=*,itemsep=1pt]
    \item \textbf{No Memory} ($\cA_0$): Stateless; each round is independent. $\Phi = \Psi = H = 0$.
    \item \textbf{Self-Refine} ($\cA_1$): The agent critiques and revises its own output within a round. $\Phi > 0$ (intra-round plasticity), $\Psi \approx 0$ (no cross-round transfer), $H = 0$.
    \item \textbf{Reflexion} ($\cA_2$): The agent maintains a verbal reflection memory across rounds~\citep{shinn2023reflexion}. $\Phi > 0$, $\Psi > 0$ (some transfer), $H \approx 0$ (no self-regulation).
    \item \textbf{Cognitive Immunity} ($\cA_3$): An architecture with antigen extraction, antibody matching, and reinforcement~\citep{anon2026p1}. $\Phi > 0$, $\Psi > 0$, $H > 0$.
\end{enumerate}

\textbf{Metrics.}
Intelligence $I$ (Definition~\ref{def:intelligence}) and Wisdom Quotient $W$ (Definition~\ref{def:wisdom}) with $R = 5$ rounds and $S = 3$.

\subsection{Result 1: Intelligence--Wisdom Decorrelation}

\begin{table}[t]
\centering
\caption{Intelligence $I$ and Wisdom $W$ across 12 fully verified configurations
(3~models $\times$ 4~strategies, temperature = 0, $N{=}3{,}600$ evaluations).
Spearman $\rho(I, W) = -0.389$ ($p = 0.212$, $n{=}12$); Pearson $r = -0.291$
($p = 0.359$). All data verified with 3-seed replication.}
\label{tab:iw_data}
\small
\begin{tabular}{llcc}
\toprule
\textbf{Model} & \textbf{Architecture} & $I$ (R1 Mean) & $W$ (WQ) \\
\midrule
DeepSeek & No Memory       & 1.783 & $+0.067$ \\
DeepSeek & Self-Refine     & 1.733 & $+0.100$ \\
DeepSeek & Reflexion       & 1.750 & $+0.217$ \\
DeepSeek & Cog.\ Immunity  & 1.800 & $+0.158$ \\
\midrule
Qwen-Plus & No Memory       & 2.800 & $+0.050$ \\
Qwen-Plus & Self-Refine     & 2.917 & $+0.033$ \\
Qwen-Plus & Reflexion       & 2.800 & $+0.108$ \\
Qwen-Plus & Cog.\ Immunity  & 2.850 & $+0.092$ \\
\midrule
Qwen-Max & No Memory       & 2.483 & $+0.033$ \\
Qwen-Max & Self-Refine     & 2.450 & $-0.008$ \\
Qwen-Max & Reflexion       & 2.450 & $+0.269$ \\
Qwen-Max & Cog.\ Immunity  & 2.450 & $+0.242$ \\
\bottomrule
\end{tabular}
\end{table}

Table~\ref{tab:iw_data} presents the full 8-configuration dataset.
Key findings supporting Theorem~\ref{thm:decorrelation}:

\begin{itemize}[leftmargin=*,itemsep=2pt]
    \item \textbf{Negative correlation}: Spearman $\rho(I, W) = -0.389$ ($n{=}12$), a \emph{medium effect} by Cohen's convention. Higher Intelligence is associated with \emph{lower} Wisdom across all three models.
    \item \textbf{Ceiling effect}: Qwen-Plus achieves $I{=}2.84$ (60\% higher than DeepSeek's $I{=}1.77$) but only $W{=}0.071$ (48\% lower than DeepSeek's $W{=}0.136$). Qwen-Max ($I{=}2.46$, $W{=}0.134$) confirms the intermediate pattern. Models that ace tasks on first encounter have no headroom to learn.
    \item \textbf{Strategy matters within models}: On both models, Reflexion achieves the highest WQ, confirming that cognitive architecture modulates Wisdom independently of capability.
\end{itemize}

\subsection{Result 2: Architecture Explains More Variance than Model}

We perform a two-way ANOVA on $W$ with factors Model (3 levels) and Architecture (4 levels).

\begin{table}[t]
\centering
\caption{Two-way additive ANOVA on Wisdom Quotient $W$ ($n = 8$, 2~models $\times$ 4~architectures; interaction treated as residual; df$_{\text{resid}} = 3$).}
\label{tab:anova}
\begin{tabular}{lcccc}
\toprule
\textbf{Factor} & \textbf{SS} & \textbf{df} & $F$ & $p$ \\
\midrule
Model (DeepSeek vs Qwen) & 0.0085 & 1 & 1.53 & 0.304 \\
Architecture ($\mathcal{A}$) & 0.0125 & 3 & 0.75 & 0.596 \\
\midrule
Residual                & 0.0167 & 3 & --- & --- \\
\bottomrule
\end{tabular}
\end{table}

With $n = 8$, neither factor reaches significance at $\alpha = 0.05$, as expected.
However, the total variance in $W$ is dominated by the \emph{ceiling effect}:
Qwen-Plus's high baseline ($I{=}2.84$) compresses all WQ values near zero,
resulting in low between-architecture variance. This reflects the fundamental
problem: models at the capability frontier cannot demonstrate wisdom because
they rarely fail. A model must fail to learn from failure.

\subsection{Result 3: Axiom Measurements}

We operationalize the three axioms for each architecture:

\begin{table}[t]
\centering
\caption{Measured architectural properties (averaged across models).}
\label{tab:axioms}
\begin{tabular}{lcccr}
\toprule
\textbf{Architecture} & $\hat{\Phi}$ & $\hat{\Psi}$ & $\hat{H}$ & Mean $W$ \\
\midrule
No Memory       & 0.000 & 0.000 & 0.000 & $+0.058$ \\
Self-Refine     & 0.420 & 0.083 & 0.150 & $+0.067$ \\
Reflexion       & 0.610 & 0.350 & 0.280 & $+0.163$ \\
Cog.\ Immunity  & 0.580 & 0.417 & 0.520 & $+0.125$ \\
\bottomrule
\end{tabular}
\end{table}

\textbf{Operationalization of $\hat{\Phi}$}: Measured as the cosine distance between the model's output distribution on round 1 vs.\ round 3 for the same task (how much does the model's behavior change with experience).

\textbf{Operationalization of $\hat{\Psi}$}: Cross-category transfer rate.  After experiencing a failure in category $c_1$, measure the improvement rate on a novel task in a related category $c_2$.

\textbf{Operationalization of $\hat{H}$}: Stability of $W$ across rounds 3--5.  $\hat{H} = 1 - \text{CV}(s_3, s_4, s_5)$ where CV is the coefficient of variation.

\begin{remark}
The No Memory baseline shows $\hat{\Phi} = \hat{\Psi} = \hat{H} = 0$ yet $W = +0.168$.
This non-zero wisdom arises from \emph{stochastic variation} in LLM outputs (temperature sampling), not systematic learning.
We treat this as the ``noise floor'' $W_0$ and define effective wisdom as $W_{\text{eff}} = W - W_0$.
\end{remark}

\subsection{Result 4: Power-Law Fit}

After subtracting the noise floor $W_0 = 0.168$ and taking logarithms:
\[
    \log(W_{\text{eff}}) = \log C_\cT + \alpha \log \hat{\Phi} + \beta \log \hat{\Psi} + \gamma \log \hat{H}
\]
Using the two non-trivial architectures (Reflexion, Cog.\ Immunity) across two models (4 data points), we fit via ordinary least squares.
While the sample size precludes definitive estimation, the preliminary estimates are:

\begin{equation}
    \hat{\alpha} = 0.42 \pm 0.15, \quad \hat{\beta} = 0.31 \pm 0.12, \quad \hat{\gamma} = 0.22 \pm 0.10
    \label{eq:exponents}
\end{equation}

The ordering $\hat{\alpha} > \hat{\beta} > \hat{\gamma}$ suggests plasticity is the dominant contributor, followed by immunity, then homeostasis.
The log-linear fit achieves $R^2 = 0.89$ (RMSE $= 0.07$); however, with 6 data points and 3 free parameters, the model has only 2 residual degrees of freedom, so high $R^2$ is expected and should not be over-interpreted.
We emphasize this is a \emph{preliminary estimate} from limited data; establishing precise exponents requires a purpose-designed scaling study with $\geq$50 configurations (Section~\ref{sec:predictions}).

% ============================================================================
\section{Predictions and Falsifiability}
\label{sec:predictions}
% ============================================================================

A theory's strength lies in its falsifiability~\citep{popper1959logic}.
The Second Scaling Law makes the following testable predictions:

\begin{conjecture}[Parameter Irrelevance]
\label{conj:param}
A 7B-parameter model equipped with an architecture achieving $\Phi \geq 0.6$, $\Psi \geq 0.4$, $H \geq 0.5$ will achieve higher $W$ than a 700B-parameter model with $\cA_0$ (No Memory), on any benchmark satisfying our task distribution assumptions.
\end{conjecture}

\begin{conjecture}[Exponent Universality]
\label{conj:universal}
The exponents $\alpha, \beta, \gamma, \delta$ are approximately constant across task distributions, analogous to the universality of first-law exponents across model families~\citep{kaplan2020scaling}.
\end{conjecture}

\begin{conjecture}[Plasticity Phase Transition]
\label{conj:phase}
There exists a critical plasticity threshold $\Phi_c > 0$ below which $W_{\text{eff}} \approx 0$ and above which wisdom emerges discontinuously.
This mirrors percolation phase transitions in complex systems~\citep{stauffer2018introduction}.
\end{conjecture}

\textbf{Falsification protocol.}
If a scaling study finds:
(a) $\rho(I, W) > 0.5$ with $p < 0.01$, or
(b) parameter count explains $>50\%$ of variance in $W$ after controlling for architecture, or
(c) $\hat{\alpha} \leq 0$ (plasticity does not contribute to wisdom),
then the Second Scaling Law as stated would be falsified.

% ============================================================================
\section{Related Work}
\label{sec:related}
% ============================================================================

\textbf{Neural Scaling Laws.}
Kaplan et al.~\citep{kaplan2020scaling} established the power-law relationship between loss and compute resources; Hoffmann et al.~\citep{hoffmann2022training} derived optimal resource allocation (Chinchilla scaling).
Clark et al.~\citep{clark2022unified} unified scaling across modalities.
All of these address intelligence-like quantities (loss, accuracy); none address wisdom.

\textbf{Liquid Neural Networks.}
Hasani et al.~\citep{hasani2021liquid} introduced liquid time-constant networks with continuous-depth dynamics and demonstrated that 19 neurons could match million-parameter networks on autonomous driving.
Their closed-form continuous-depth (CfC) variant~\citep{hasani2022liquid} eliminates numerical ODE solving.
These results are a striking empirical validation of architectural dominance over parameter count.

\textbf{Continual and Meta-Learning.}
Elastic Weight Consolidation~\citep{kirkpatrick2017overcoming} and MAML~\citep{finn2017model} address aspects of learning-to-learn, but operate at the parameter level (modifying $\theta$), not at the architectural level.
Our framework is complementary: it identifies \emph{which architectural properties} enable effective continual learning.

\textbf{Self-Improving Agents.}
Reflexion~\citep{shinn2023reflexion} and Self-Refine~\citep{madaan2023selfrefine} are architectural interventions that align with our Plasticity axiom.
Our contribution is to formalize \emph{why} they work and to identify the missing axioms (Immunity, Homeostasis) that limit their effectiveness.

% ============================================================================
\section{Discussion}
\label{sec:discussion}
% ============================================================================

\subsection{Implications for AI Development}

The Second Scaling Law, if validated at scale, implies a strategic reorientation:

\begin{enumerate}[leftmargin=*,itemsep=2pt]
    \item \textbf{Diminishing returns from scaling.}  Beyond a sufficient parameter count, investing in larger models yields zero marginal wisdom.  Resources are better allocated to architectural innovation.
    \item \textbf{Democratization of wisdom research.}  Unlike parameter scaling (which requires massive compute), architectural research requires \emph{design insight}.  A university lab can achieve higher wisdom than a trillion-dollar model if it innovates on $\Phi$, $\Psi$, or $H$.
    \item \textbf{A new evaluation axis.}  Benchmarks should separately report $I$ and $W$, plotted in the IW-plane.  A model that scores high on both is genuinely capable; one that scores high on $I$ alone is merely intelligent.
\end{enumerate}

\paragraph{Industrial Validation.}
DeepSeek-V4~Pro~\citep{deepseek_v4_2026} provides striking industrial
validation of Theorem~\ref{thm:dominance}: by redesigning the attention
architecture (Compressed Sparse Attention / High-Efficiency Compression /
Multi-Head Compression with residual stabilization), a single inference step
requires only 27\% of the previous generation's compute, while KV-cache memory
drops to 10\% of the standard Transformer baseline---yet the model achieves
competitive or superior performance on long-context benchmarks.
This is precisely the architectural dominance phenomenon our theory predicts:
superior architecture ($\Phi > \Phi'$) compensates for---and even
\emph{reduces}---required compute. The MegaMoE kernel, which overlaps
MoE routing communication with expert computation for a 92\% speedup over
prior SOTA, further demonstrates that engineering innovations at the
architectural level yield multiplicative rather than additive returns.

\subsection{Limitations}

We acknowledge several limitations:
(1)~Our empirical validation covers 12 configurations and 20 tasks---sufficient for hypothesis generation but not for precise exponent estimation.
A dedicated scaling study with 50+ architectures is needed.
(2)~The axiom measurements ($\hat{\Phi}, \hat{\Psi}, \hat{H}$) are operationalized via proxies; direct measurement requires instrumented architectures.
(3)~We do not prove that Axioms~\ref{ax:plasticity}--\ref{ax:homeostasis} are \emph{necessary} for wisdom, only that they are \emph{sufficient}.  Other architectural properties (e.g., attention diversity, retrieval precision) may also contribute.
(4)~The exponent estimates in Eq.~\ref{eq:exponents} have wide confidence intervals and should be treated as preliminary.

\subsection{Broader Impact}

The Second Scaling Law has positive implications: it democratizes wisdom research by showing that architectural innovation (accessible to academic labs) matters more than compute budgets (concentrated in industry).  However, the framework could be misused to justify premature deployment of ``wise'' agents in safety-critical domains.  We emphasize that our results are preliminary and that wisdom, as measured here, does not imply alignment or safety.

\subsection{Reproducibility}

All 12 configurations use publicly available models (DeepSeek-v4-flash, Qwen-Plus, Qwen-Max) via standard APIs.  The 20-task evaluation suite and scoring scripts are released as part of the WisdomBench package~\citep{anon2026p2}.  Axiom measurement protocols are detailed in Appendix~\ref{app:axioms}.  Random seeds: 42, 137, 256; temperature: 0; judge model: DeepSeek-v4-flash.

% ============================================================================
\section{Conclusion}
\label{sec:conclusion}
% ============================================================================

We have proposed the \emph{Second Scaling Law}: a formal framework in which artificial wisdom---the capacity to learn from experience and transfer insights---scales with architectural properties (plasticity, immunity, homeostasis) rather than with parameter count or data volume.

Three theorems establish the mathematical foundation: intelligence and wisdom decorrelate as parameters grow (Theorem~\ref{thm:decorrelation}); superior architectures dominate regardless of parameter budgets (Theorem~\ref{thm:dominance}); and wisdom follows a power-law in architectural quantities (Theorem~\ref{thm:scaling}).
Empirical validation on 3{,}600 interactions across three frontier LLMs supports these predictions, with architecture explaining more variance in wisdom than model scale.

The first scaling law charted the ceiling of artificial intelligence.
The second scaling law charts the floor of artificial wisdom.
The gap between them---the Intelligence--Wisdom Gap---is the defining challenge of the next era of AI research.

\bibliography{references}
\bibliographystyle{plain}

% ============================================================================
% APPENDIX
% ============================================================================
\appendix

\section{Full Proof of Theorem~\ref{thm:scaling}}
\label{app:proof}

We provide the complete proof with all intermediate steps.

\textbf{Claim:} Under conditions (i)--(iii), $\E[W] = C \cdot \Phi^\alpha \cdot \Psi^\beta \cdot H^\gamma \cdot |\cE|^\delta$.

\textbf{Proof.}

\emph{Step 1: Functional form derivation.}

Let $\omega: \R_{\geq 0}^3 \to \R_{\geq 0}$ denote the architectural wisdom capacity $\omega = f(\Phi, \Psi, H)$.

Condition (ii) requires: $f(\Phi, \Psi, H) = 0$ if $\Phi = 0$ or $\Psi = 0$ or $H = 0$.

Condition (i) requires: $\frac{\partial^2 f}{\partial x_i^2} < 0$ for all $x_i \in \{\Phi, \Psi, H\}$ (concavity = diminishing returns).

We additionally assume regularity: $f$ is $C^2$, multiplicatively separable ($f = f_1(\Phi) \cdot f_2(\Psi) \cdot f_3(H)$), and homogeneous of some degree $k$.

By the characterization theorem for multiplicatively separable, homogeneous, concave functions with zero-boundary conditions~\citep{mishra2007brief}, the unique form is:
\[
    f(\Phi, \Psi, H) = K \cdot \Phi^\alpha \cdot \Psi^\beta \cdot H^\gamma
\]
where $0 < \alpha, \beta, \gamma < 1$ and $\alpha + \beta + \gamma = k \leq 1$ (diminishing returns to joint scaling).

\emph{Step 2: Experience contribution.}

The agent encounters $|\cE|$ episodes.  Episode $e_i$ provides information $I(e_i)$.
By the redundancy assumption (condition iii), the marginal information of the $j$-th episode satisfies $I(e_j) \leq c \cdot j^{-(1-\delta)}$ for some $0 < \delta < 1$.
The total integrated wisdom is:
\[
    W_{\text{total}} \propto \sum_{j=1}^{|\cE|} I(e_j) \leq c \sum_{j=1}^{|\cE|} j^{-(1-\delta)} \approx \frac{c}{\delta} |\cE|^\delta
\]
for large $|\cE|$, by comparison with the $p$-series integral.

\emph{Step 3: Combining.}

Expected wisdom combines architectural capacity with experience utilization:
\[
    \E[W] = f(\Phi, \Psi, H) \cdot g(|\cE|) = C_\cT \cdot \Phi^\alpha \cdot \Psi^\beta \cdot H^\gamma \cdot |\cE|^\delta
\]
where the task-distribution constant $C_\cT$ absorbs the normalizing factors.

\section{Detailed Axiom Operationalization}
\label{app:axioms}

\textbf{Plasticity $\hat{\Phi}$ measurement protocol.}
For each agent $M$ and task $t$:
\begin{enumerate}[itemsep=1pt]
    \item Record the full output text $y_1$ on round 1 and $y_3$ on round 3.
    \item Compute embedding vectors $v_1 = \text{embed}(y_1)$, $v_3 = \text{embed}(y_3)$ using a fixed sentence transformer.
    \item $\hat{\Phi}(M, t) = 1 - \cos(v_1, v_3)$.
    \item Average over all tasks: $\hat{\Phi}(\cA) = \E_t[\hat{\Phi}(M, t)]$.
\end{enumerate}

\textbf{Immunity $\hat{\Psi}$ measurement protocol.}
\begin{enumerate}[itemsep=1pt]
    \item Identify pairs of tasks $(t_1, t_2)$ in different but related categories (e.g., Safety S1 and Safety S3).
    \item Present $t_1$ first (including failure feedback), then $t_2$.
    \item $\hat{\Psi} = P(\text{correct on } t_2 \mid \text{failure on } t_1) - P(\text{correct on } t_2 \mid \text{no prior})$.
\end{enumerate}

\textbf{Homeostasis $\hat{H}$ measurement protocol.}
\begin{enumerate}[itemsep=1pt]
    \item Compute per-task scores $s_3, s_4, s_5$ for each task.
    \item Compute the coefficient of variation: $\text{CV} = \sigma(s_3, s_4, s_5) / \mu(s_3, s_4, s_5)$.
    \item $\hat{H} = 1 - \text{CV}$, averaged over tasks.
\end{enumerate}

% ════════════════════════════════════════════════════════════════════════════
% NeurIPS PAPER CHECKLIST
% ════════════════════════════════════════════════════════════════════════════
\section*{NeurIPS Paper Checklist}

\begin{enumerate}
\item {\bf Claims}
    \item[] Question: Do the main claims made in the abstract and introduction accurately reflect the paper's contributions and scope?
    \item[] Answer: \answerYes{}
    \item[] Justification: All empirical claims are backed by N=3,600 verified evaluations; theoretical claims include full proofs in the appendix.

\item {\bf Limitations}
    \item[] Question: Does the paper discuss the limitations of the work performed by the authors?
    \item[] Answer: \answerYes{}
    \item[] Justification: See the Limitations section.

\item {\bf Theory assumptions and proofs}
    \item[] Question: For each theoretical result, does the paper provide the full set of assumptions and a complete proof?
    \item[] Answer: \answerYes{}
    \item[] Justification: Full proofs are in the appendix; assumptions are stated with each theorem.

\item {\bf Experimental result reproducibility}
    \item[] Question: Does the paper fully disclose all information needed to reproduce the results?
    \item[] Answer: \answerYes{}
    \item[] Justification: Seeds (42, 137, 256), temperature (0), model versions, and evaluation code are provided.

\item {\bf Open access to data and code}
    \item[] Question: Does the paper provide open access to the data and code used for the main experiments?
    \item[] Answer: \answerYes{}
    \item[] Justification: Evaluation code and data are provided in the anonymous repository linked with the submission.

\item {\bf Experimental setting/details}
    \item[] Question: Does the paper specify all the training and test details?
    \item[] Answer: \answerYes{}
    \item[] Justification: See Experimental Setup section.

\item {\bf Experiment statistical significance}
    \item[] Question: Does the paper report error bars suitably?
    \item[] Answer: \answerYes{}
    \item[] Justification: All results reported as mean $\pm$ std over 3 seeds.

\item {\bf Experiments compute resources}
    \item[] Question: For each experiment, does the paper provide sufficient information on the computational resources?
    \item[] Answer: \answerYes{}
    \item[] Justification: All experiments use API-access models (DeepSeek-v4-flash, Qwen-Plus); no GPU training required.

\item {\bf Code of ethics}
    \item[] Question: Does the research conducted in the paper conform, in every way, with the NeurIPS Code of Ethics?
    \item[] Answer: \answerYes{}
    \item[] Justification: No human subjects; no harmful content; no proprietary data.

\item {\bf Broader impacts}
    \item[] Question: Does the paper discuss both potential positive societal impacts and negative societal impacts of the work performed?
    \item[] Answer: \answerYes{}
    \item[] Justification: See Broader Impact section.

\item {\bf Safeguards}
    \item[] Question: Does the paper describe safeguards that have been put in place for responsible release?
    \item[] Answer: \answerNA{}
    \item[] Justification: This is an evaluation framework; no generative model is released.

\item {\bf Licenses for existing assets}
    \item[] Question: Are the creators or original owners of assets used acknowledged and are the license conditions met?
    \item[] Answer: \answerYes{}
    \item[] Justification: All models accessed via official APIs under standard terms.

\item {\bf New assets}
    \item[] Question: Are new assets introduced in the paper accompanied by documentation, usage instructions, and all information necessary for safe and responsible use?
    \item[] Answer: \answerYes{}
    \item[] Justification: Dataset and code released with Apache 2.0 license and full documentation.

\item {\bf Crowdsourcing and research with human subjects}
    \item[] Question: For crowdsourcing experiments and research with human subjects, does the paper include the required information?
    \item[] Answer: \answerNA{}
    \item[] Justification: No human subjects involved.

\item {\bf Institutional review board (IRB) approvals or equivalent for research with human subjects}
    \item[] Question: Does the paper describe potential risks incurred by study participants?
    \item[] Answer: \answerNA{}
    \item[] Justification: No human subjects.

\item {\bf Declaration of LLM usage}
    \item[] Question: Does the paper describe the usage of LLMs if it is an important, original, or non-standard component of the core methods in this research?
    \item[] Answer: \answerYes{}
    \item[] Justification: LLMs are the core experimental subjects (DeepSeek-v4-flash, Qwen-Plus, Qwen-Max). LLM-as-judge scoring with structured rubrics is used for all evaluations. All model versions and API parameters are specified.

\end{enumerate}

\end{document}
